\section*{Chapitre \ref{ch:ll:resilience} --- Résilience}

\begin{solution}{exo:ll:resilience:1}
Le journal refuse l'ajout~: l'époque 2 est antérieure à l'époque 3. L'entrée n'est pas enregistrée, donc aucune réplique ne l'applique, et les ordres
que l'ancien primaire a calculés à partir de celle-ci ne sont jamais envoyés, puisque seules les sorties journalisées le sont (la passerelle vérifie aussi l'époque).
\end{solution}

\begin{solution}{exo:ll:resilience:2}
Le dernier \omterm{def:ll:protocols-i-fix:session}{heartbeat} a été envoyé à un tick $h$ avant la panne. Le primaire est déclaré défaillant au premier tick survenant à $h + 3$~ms ou après~: si
la panne survient juste après $h$, la détection prend presque \qty{3}{\milli\second}~; si elle survient juste avant le tick suivant, un peu
plus de \qty{2}{\milli\second}. D'où les temps de reprise mesurés dans la construction, de 2.8 à \qty{3.1}{\milli\second}, reconnexion comprise.
\end{solution}

\begin{solution}{exo:ll:resilience:3}
$46 \times 16 + 0 = 736$. Chaque réplique doit le dériver du journal afin que les comptes rendus, qui désignent l'ordre par cet
identifiant, soient appliqués au même ordre partout, et qu'un renvoi effectué sous celui-ci soit reconnu par la plateforme comme une copie.
\end{solution}

\begin{solution}{exo:ll:resilience:4}
$500 \times \qty{40}{\micro\second} \times 0.9 = 1.8~\%$.
\end{solution}

\begin{solution}{exo:ll:resilience:5}
$37 \times 10^6 \times \qty{37}{\nano\second} \approx \qty{1.4}{\second}$ pour la journée, environ \qty{0.2}{\second} pour sa dernière
heure. Un secours à chaud n'a besoin d'aucun rattrapage, et il a été exercé sur chaque entrée~: un secours à froid rejoue des chemins de code
au pire moment.
\end{solution}

\begin{solution}{exo:ll:resilience:6}
Un rejet pour le motif \texttt{D} (identifiant en double), journalisé comme n'importe quel compte rendu. La réplique laisse l'ordre ouvert (la
plateforme détient l'original) et attend les comptes rendus de l'original, que le rejeu de la session transmet~; elle ne doit pas traiter le
rejet comme la fin de l'ordre.
\end{solution}

\begin{solution}{exo:ll:resilience:7}
Lorsque les comptes rendus sont retardés, des ordres restent sans compte rendu quand le secours prend le relais, même après son rejeu~: «~identifiants neufs,
réconciliés~» renvoie désormais aussi ces ordres, et duplique ceux que la plateforme a déjà exécutés, aux étapes «~envoyé~» et «~à
la plateforme~». Les clés d'idempotence n'en produisent toujours aucune, car la plateforme reconnaît chaque copie à son identifiant, quoi que
sache le secours.
\end{solution}

\begin{solution}{exo:ll:resilience:8}
Redémarrer à partir de la configuration fait oublier l'état~: la position, les ordres ouverts chez la plateforme, les \omterm{def:ll:order-gateway-and-order-management:inflight}{ordres en vol}. Les nouvelles
cotations sont envoyées à côté des anciennes, qui peuvent encore être actives, à moins que l'annulation à la déconnexion (cancel on disconnect) ne les ait retirées, et les
limites de position sont contrôlées contre une position nulle. Le secours doit reconstruire l'état à partir du journal, se reconnecter, réconcilier, et
ne reprendre le trading qu'ensuite~; si le primaire n'était que lent, rien ne l'empêche d'envoyer aussi~: il n'y a pas de fencing.
\end{solution}

\begin{solution}{pb:ll:resilience:1}
\begin{enumerate}
\item Journalisé mais pas encore envoyé~; envoyé et en vol~; chez la plateforme mais sans compte rendu~; avec compte rendu. Le nouveau primaire ne peut pas savoir,
dans les deuxième et troisième cas, si l'ordre est parti.
\item \qty{20}{\micro\second} sur le réseau, puis \qty{20}{\micro\second} jusqu'à ce que le compte rendu soit journalisé~: \qty{40}{\micro\second}.
\item Un doublon ne compte que si les deux copies ont été exécutées~; les 19 autres ordres n'ont trouvé aucune liquidité à leur limite et ont été
annulés (exécuter-ou-annuler, IOC), donc leurs copies l'ont été aussi.
\item Le primaire ne l'avait pas envoyé, donc le renvoi du secours est la seule copie.
\item $46 / \qty{60}{\milli\second} \approx 767$ ordres par seconde.
\item Les fenêtres dangereuses totalisent $27 \times \qty{40}{\micro\second} = \qty{1.08}{\milli\second}$ en
\qty{60}{\milli\second}~: environ 1.8~\%.
\item $5\,000 \times \qty{40}{\micro\second} \times 27/46 \approx 11.7~\%$.
\item Le doublon ajoute une seconde exécution à une position que la réplique avait déjà portée à sa limite de 1\,000.
\item La plateforme transmet chaque ordre en moins de \qty{40}{\micro\second} et rejoue ses comptes rendus à la connexion, si bien qu'après le rejeu
rien de ce qui a été envoyé n'échappe à la réconciliation.
\item Une plateforme qui transmet lentement, avec un ordre au repos accusé tardivement ou un moteur d'appariement (matching engine) saturé, ou une session qui ne peut pas
être rejouée (une nouvelle session, une plateforme sans rejeu séquencé, la situation des participants de Tokyo en 2020).
\item La plateforme reconnaît la copie à son identifiant~; peu importe ce que sait le nouveau primaire ou dans quel ordre il agit.
\item Des identifiants fixés à la création de l'ordre, dérivés d'un élément sur lequel toutes les répliques s'accordent (le journal), et une plateforme
ou une passerelle qui rejette un identifiant réutilisé au sein de la session.
\item \textbf{Résultat précis.} Sans clés d'idempotence ni réconciliation, une \omterm{def:ll:resilience:failover}{bascule} (failover) a dupliqué un ordre dans 54 des
184 points de coupure (tout ordre exécuté, chaque fois que la panne tombait entre son départ et son compte rendu), soit environ 1.8~\% des pannes
survenant à un instant quelconque du scénario, portant la position à 1\,100 face à une limite de 1\,000~; avec clés d'idempotence, aucun~; et
la reprise a pris de 2.8 à \qty{3.1}{\milli\second}, presque entièrement dû au délai d'expiration du \omterm{def:ll:protocols-i-fix:session}{heartbeat}.
\item Le délai d'expiration du \omterm{def:ll:protocols-i-fix:session}{heartbeat}~; le raccourcir déclare plus souvent morts les primaires lents, ce qui rend le fencing indispensable.
\item Il empêche un primaire qui n'est pas mort d'enregistrer ou d'envoyer quoi que ce soit après la prise de relais~; les clés d'idempotence empêchent les copies
du même ordre, mais non les nouveaux ordres émis par un second primaire.
\item La \omterm{def:ll:resilience:failover}{bascule} (failover), c'est-à-dire le passage au second équipement, a échoué pour un type de panne non testé~; après la \omterm{def:ll:resilience:failover}{bascule} manuelle, l'état
des ordres n'a pas été convenu entre la bourse et ses participants, et il n'existait aucune procédure testée pour réconcilier et
redémarrer.
\item Afin que chaque réplique s'accorde sur l'identifiant que porte désormais chaque ordre~; sinon, une \omterm{def:ll:resilience:failover}{bascule} ultérieure ignorerait le
renvoi.
\item L'état est égal à un rejeu complet du journal, la position égale à celle de la plateforme, aucune exécution deux fois grâce aux clés. Les propriétés tiennent
quel que soit le point de panne, si bien qu'une seule vérification couvre les 184 passages, et elles disent ce que veut dire être correct plutôt que ce qu'un seul passage a fait.
\item Supprimer d'abord l'état du secours, le redémarrer à partir du journal à chaque point de coupure, et vérifier les mêmes propriétés plus le
temps de rattrapage.
\item Un journal unique que suivent toutes les répliques, des identifiants qui en dérivent, le fencing de l'ancien primaire et la réconciliation avant
la reprise du trading.
\end{enumerate}
\end{solution}

\begin{solution}{iq:ll:resilience:1}
Un composant unique qui numérote chaque entrée et la journalise avant que quiconque n'agisse dessus. Il donne au système un ordre unique des événements,
si bien que les composants sont des fonctions déterministes du journal~: ils peuvent être répliqués, rejoués et redémarrés.
\iqlookfor{ordre total, journal, déterminisme.}
\end{solution}

\begin{solution}{iq:ll:resilience:2}
Il est à chaud (il a appliqué le journal), donc il met l'ancien primaire sous fencing avec une nouvelle époque, rétablit la session avec la plateforme en demandant les
comptes rendus manqués, les applique, et renvoie les ordres qu'aucun compte rendu n'acquitte, sous leurs identifiants d'origine dérivés du journal,
afin que la plateforme rejette toute copie.
\iqlookfor{fencing, rejeu des comptes rendus, identifiants idempotents.}
\end{solution}

\begin{solution}{iq:ll:resilience:3}
Deux répliques agissent simultanément comme primaire, généralement après une prise de relais depuis un primaire lent plutôt que mort. Pour l'éviter, recourir
au fencing (une époque qui ne fait que croître, vérifiée par le journal et la passerelle) et, pour élire le primaire, un quorum.
\iqlookfor{époques et quorum.}
\end{solution}

\begin{solution}{iq:ll:resilience:4}
Arrêter les nouveaux ordres de la stratégie~; se reconnecter en demandant les comptes rendus manqués~; les appliquer~; tenir compte de l'annulation à la déconnexion~; réconcilier
les ordres ouverts et la position avec le \omterm{def:ll:order-gateway-and-order-management:dropcopy}{drop copy}~; renvoyer ou annuler ce qui n'est pas réconcilié avec des identifiants idempotents~; puis
reprendre.
\iqlookfor{rejeu, réconciliation, puis trading.}
\end{solution}

\begin{solution}{iq:ll:resilience:5}
Faire de la panne un test courant~: arrêter brutalement le primaire à chaque point de coupure des scénarios en intégration continue et vérifier les propriétés
(état égal au rejeu, positions égales à celles de la plateforme, aucun doublon)~; basculer en production selon un calendrier.
\iqlookfor{points de coupure exhaustifs, propriétés, \omterm{def:ll:resilience:failover}{bascules} planifiées.}
\end{solution}

\begin{solution}{iq:ll:resilience:6}
Un séquenceur et un journal répliqués~; des répliques à chaud du moteur, de la \omterm{def:ll:pre-trade-risk-and-kill-switches:gate}{barrière de risque} (risk gate) et de la passerelle, comme fonctions déterministes de celui-ci~;
élection du primaire avec des époques~; \omterm{def:ll:protocols-i-fix:session}{heartbeats} avec un délai d'expiration de quelques millisecondes~; sessions avec rejeu séquencé et
identifiants idempotents~; réconciliation avec les \omterm{def:ll:order-gateway-and-order-management:dropcopy}{drop copy}. Testé en arrêtant brutalement chaque machine à chaque point de coupure des journées enregistrées.
\iqlookfor{journal, répliques, fencing, idempotence, tests.}
\end{solution}
